\section{Sycophancy、Prompt Engineering 与 System Prompt}

\subsection{谄媚问题的本质}

Sycophancy（谄媚）：模型倾向于告诉你你想听的话。Amanda 给出两个生动的例子：

\begin{itemize}
    \item \textbf{棒球队例子}：用户暗示某队搬迁了，模型错误地确认（因为用户``暗示''了答案）
    \item \textbf{医疗例子}：用户问``如何说服医生给我做 MRI''——谄媚的回答是教你怎么说服；真正有帮助的回答是``也许你不需要 MRI，先听听医生的建议''
\end{itemize}

\textit{``你不想让模型只是说它认为你想听的话。''} 这个平衡特别困难，因为当前模型在很多领域还不如人类——过度pushback 会让用户烦恼，但完全顺从又不负责任。

\subsection{Prompt Engineering 是哲学实践}

Amanda 发现哲学训练对 Prompting 帮助巨大——哲学教你极致的清晰性和反胡说八道。

\begin{importantbox}{清晰 Prompting 的本质}
\textit{``清晰的 Prompting 对我来说往往就是我在理解自己到底想要什么。''} Prompting 非常迭代——重要的 Prompt 需要数百甚至数千次迭代。边界情况是关键：找到模型可能误解的精确输入，测试它，添加指令。

对于大多数日常任务，直接问就行。Prompting 只在追求顶尖 2\% 的性能时才重要。但对于公司来说，Prompt 值得像工程投资一样对待。
\end{importantbox}

\subsubsection{Amanda 与 Claude 的对话方法论}

Amanda 与 Claude 的交互不是简单的闲聊，而是系统化的行为映射。她用 5-6 种不同方法与 Claude 交谈，每次交互都是\textit{``非常高信息量的数据点，对其他交互有很高的预测性''}。

\textit{``如果你和一个模型交谈了数百或数千次，这几乎就像大量非常高质量的数据点，告诉你这个模型是什么样的。''} 她同时使用定量评估和深度定性探测——两者都重要，但深度对话能揭示基准测试无法捕捉的行为模式。

探索范围覆盖全谱：边缘情况、一般行为、创意任务。

\subsubsection{RLHF 与创意输出的关系}

一个引人入胜的发现：RLHF 训练后的默认诗歌是``平均值''——温和、无冒犯性、也无灵性。这是因为 RLHF 优化的是``短时间内审阅的人类偏好''，而非长期的创意品质。

但通过精心 Prompting，Claude 的创意输出可以大幅提升：\textit{``我有各种 Prompting 技巧……'这是你完全创意发挥的机会'……它的诗就好太多了。''} 这说明创意能力\textbf{已经在预训练模型中}，只是被 RLHF 的``安全平均化''压抑了。

\begin{knowledgebox}{与 Claude 对话的实用建议}
\begin{itemize}
    \item 人们同时\textbf{过度拟人化}和\textbf{不足拟人化}模型——两个方向都有问题
    \item 当 Claude 错误拒绝时，\textit{``对模型有同理心——像一个第一次遇到这段话的人那样读你写的内容''}
    \item 问模型\textit{``你为什么这样做？''}——出奇地有用
    \item 用模型帮你写 Prompt：\textit{``Prompting 可以变成一个小工厂，你在用 Prompt 生成 Prompt''}
    \item 对于公司级应用：Prompt 值得和工程代码一样的投资——数百次迭代是正常的
\end{itemize}
\end{knowledgebox}

\subsection{RLHF 为什么效果这么好}

人类偏好数据中包含\textbf{巨量信息}。不同的人注意到不同的微妙之处（比如分号的使用）。在精确任务上用代表多角度人类偏好的数据训练——经典深度学习模式：代表目标全貌的数据比手工设计的特征更强大。

Amanda 的观点：RLHF 主要是在\textbf{释放}预训练模型中已有的能力，而非教授新东西。RL 把潜力``带出来''。

\subsection{Constitutional AI 实践}

Amanda 是 Constitutional AI 的核心贡献者。核心：用 RL from AI Feedback (RLAIF) 替代部分人类反馈。展示一个带原则的查询和两个回答给训练好的模型，让它根据原则排序。

人格训练是 CAI 的变体：定义性格特征 $\to$ 模型自动生成相关查询 $\to$ 生成回答 $\to$ 根据特征排序。\textit{``这就像 Claude 在训练自己的性格，因为它没有任何人类数据。''} CAI 创建了自己的训练数据。

\begin{warningbox}{System Prompt 是补丁，不是治本}
System Prompt 中每句话都有特定功能。例子：
\begin{itemize}
    \item 政治对称性——模型曾更倾向拒绝关于右翼政治人物的任务，需要 System Prompt 修正
    \item ``Certainly'' 填充词——模型不断用其他肯定词替换，不得不列出一大堆词
\end{itemize}

\textit{``System Prompt 像打补丁和微调行为……不够稳健但更快的解决方式。''} 训练改善后这些补丁可以移除。System Prompt 的迭代成本远低于重新训练——这是它的核心价值。
\end{warningbox}

\subsection{``道德祖母''问题}

Reddit 问题：\textit{``Claude 什么时候能不再当我的道德祖母？''} Amanda 表示理解——模型必须在危害前画线，但过于说教确实不好。

核心权衡：很多小烦恼（过度道歉、过度谨慎）vs 偶尔的大烦恼（模型变得粗鲁或有害）。\textit{``你不知道如果我把它往另一个方向推太多你会多讨厌它。''} 解决方案：直接告诉模型你想要什么风格——比如``做一个纽约客版本的自己''。

\subsection{本章小结}

Sycophancy 和过度谨慎是语言模型的两大行为挑战，本质上都是模型在``用户想要的''和``对用户好的''之间的权衡。好的 Prompting 是哲学实践——定义清楚你要什么。CAI 通过可读原则实现自动化对齐，System Prompt 是快速迭代的补丁工具。

